\section{本章小结}
\label{sec:homol-summary}

本章开头问：三角形框与三角形纸片有同样的顶点和边，怎样用代数量区分“有洞”与“被填上”；环带与圆周不同胚，怎样说明它们有同样的洞。前一个问题由单纯同调回答，后一个问题由同伦与同伦不变性回答。

同伦用乘积空间$X\times[0,1]$上的连续映射描述映射的连续变形，它是与复合相容的等价关系。把“互逆”放宽为“复合同伦于恒等映射”，得到同伦等价；与单点同伦等价的空间称为可缩空间，凸集都可缩。形变收缩把空间收回子空间并在过程中固定子空间，包含映射因而是同伦等价，环带由此与圆周同伦等价。同伦等价比同胚宽松：区间、圆盘和单点同伦等价，却两两不同胚。只有回缩而没有形变则不够，两点离散空间就是例子。

有限单纯复形用对取面封闭的顶点子集族记录点、边和各维填充，一维复形就是图。几何实现把每个单纯形换成标准单纯形，沿公共面粘合；它既是商空间，又同胚于欧氏空间中以标准基向量为顶点的凸包之并，其中两个单纯形只在公共面上相交。画在纸上的交叉不是复形中的点。

单纯同调在实数系数下取链空间$C_k(\mathcal K)\cong\mathbb R^{n_k}$，由交错取面定义边界映射，并验证$\partial\circ\partial=0$。闭链是边界映射的核，边界链是高一维边界映射的像，同调是二者之商，Betti数为$\beta_k=n_k-\operatorname{rank}\bm D_k-\operatorname{rank}\bm D_{k+1}$。$\beta_0$计数一维骨架的连通分量；三角形边界的$\beta_1=1$，填实后变为$0$；四面体表面的$\beta_2=1$。Euler示性数等于Betti数的交错和，可以核对手算。链映射诱导同调映射；比较子复形与大复形时，诱导映射能指出哪个环被填掉，单看Betti数做不到。

奇异同调以标准单纯形到空间的连续映射为基，对一切拓扑空间有定义，零维同调计数路径分支，连续映射诱导同调映射，同胚给出同构。在有限复形上，把单纯形看作线性的奇异单形，得到的比较映射在同调上是同构，本章只证明了零维情形并解释了一般机制；它使连续空间的同调可以用三角剖分的边界矩阵计算，并与三角剖分的选择无关。棱柱构造证明同伦映射诱导相同的同调映射，所以同伦等价的空间有同构的同调。由此，环带与圆周的$\beta_1$都是$1$，环带不可缩，圆盘不能回缩到边界圆周。

同调的使用有明确的方向：Betti数不同，就足以断定两个空间既不同胚，也不同伦等价；Betti数相同则什么也不能断定，圆盘与单点就是同调相同而不同胚的例子。流形上的同调在第~\ref{chap:manifold-basics}章起继续使用，从有限点云恢复同调、沿一列复形追踪诱导映射的持久同调在第~\ref{chap:manifold-recovery}章展开。下一部分转向另一类问题：不再研究连续变形与孔洞，而是在第~\ref{chap:measure}章度量集合的大小。
